\songtitle{Tiempo Real}

\tag*{2026}\tag{corporate-music}\tag{original-lyrics}\tag{spanish-lyrics}
\begin{notes}
Written by Carlos Thompson as an elevator pitch for a proposal.
\end{notes}
\begin{style}
Corporate background music with spoken word narration, A clean electric guitar plays a repetitive, palm-muted eighth-note pattern on a single note, providing a steady rhythmic pulse, A bright, acoustic piano enters with syncopated, high-register chords, The percussion consists of a light, electronic kick drum on every beat and a crisp clap on beats two and four, A subtle synth pad provides harmonic filling in the background, The tempo is 120 BPM in the key of C major, The male narration is dry with minimal reverb, delivered in a professional, informative tone
\end{style}
\begin{lyrics}
\songpart{Intro}
\annotation{palm-muted electric guitar eighth notes, electronic kick drum}

\songpart{Verse 1}
\annotation{male narration enters}\\
Cada ticket que llega a soporte toma entre dos y tres horas antes de tener una primera respuesta. \annotation{piano chords enter} Eso no es porque el problema sea complejo, es porque alguien tiene que leerlo, preguntarle a otros, entender el contexto y decidir qué hacer. Este tiempo es puro overhead. Y cuando hay reprocesos planificados, que son el caso más frecuente, el consultor invierte minutos reales, ajusta fechas, lanza el proceso, valida el cargue, pero el ticket puede demorar días enteros esperando en cola. El cliente ve días de espera, el consultor vio veinte minutos de trabajo. Además, todos los días cada consultor debe reportar manualmente el estado de sus tickets en Mantis BT. Todos los días.

\songpart{Verse 2}
Un sistema de IA agentiva elimina los tres cuellos de botella. Clasifica y enruta el ticket en segundos, gestiona la cola de reprocesos automáticamente priorizando por ventanas de tiempo disponibles y actualiza el estado del ticket sin que el consultor tenga que escribir nada al final del día.

\songpart{Outro}
El resultado: primera respuesta en minutos, no horas; reprocesos planificados sin cola de espera artificial; y el equipo libera tiempo real que hoy va en coordinación y reportes. Con ciento ochenta y siete tickets resueltos en el historial, ya tenemos los datos para entrenar el primer modelo. ¿Damos treinta minutos esta semana para revisar los números juntos?
\end{lyrics}

